\documentclass[11pt]{article}

\usepackage[margin=1in]{geometry}
\usepackage{amsmath,amssymb}
\usepackage{array}
\usepackage{booktabs}
\usepackage{enumitem}
\usepackage{hyperref}
\usepackage{longtable}
\usepackage{tabularx}
\usepackage{microtype}
\usepackage[T1]{fontenc}
\usepackage{lmodern}

\hypersetup{
    colorlinks=true,
    linkcolor=blue,
    urlcolor=blue
}

\newcommand{\MUST}{\textbf{MUST}}
\newcommand{\MUSTNOT}{\textbf{MUST NOT}}
\newcommand{\SHOULD}{\textbf{SHOULD}}
\newcommand{\SHOULDNOT}{\textbf{SHOULD NOT}}
\newcommand{\MAY}{\textbf{MAY}}
\newcolumntype{Y}{>{\raggedright\arraybackslash}X}

\title{\textit{cs-mail} Protocol Specification\\
\large Permission, Bonded Contact, and Relationship Settlement}
\author{C-SQD}
\date{Draft 0.6 --- September 2026}

\begin{document}
\maketitle

\begin{abstract}
This document defines the normative state, commands, events, settlement rules,
error conditions, and invariants of the \textit{cs-mail} protocol. The protocol
represents communication permission as a directed relationship between public
protocol identities. An unaccepted sender reserves a bond before attempting
contact. Acceptance and rejection are relationship-level decisions that atomically
settle every eligible admitted bond and fully cancel every unadmitted reservation
for that relationship. Rejection declines one
solicitation while preserving the possibility of a later attempt after backoff;
a separate recipient block prohibits future ordinary contact until removed. The
specification also distinguishes reservation from admission: an abandoned
unadmitted reservation returns in full, and the recipient decision window begins
only when the message is admitted. Repeated-attempt controls remain scoped to the
private principal and are not reset by accepting one public identity. The
specification further defines persistence reserves, expiry ordering, and the
coalescing of repeated recipient solicitations. It separates protocol settlement,
ledger operation, and external financial rails. It also defines a native
privacy profile covering end-to-end content encryption, pairwise-scoped
identifiers, compartmentalized protocol data, limited retention, and
cryptographic erasure.
\end{abstract}

\section{Status and Conformance}

This is an initial protocol draft. It is intended to make the deterministic core
implementable and testable; it does not select a wire encoding, signature suite,
database, transport, banking partner, or jurisdiction-specific compliance model.

The key words \MUST, \MUSTNOT, \SHOULD, \SHOULDNOT, and \MAY{} in this document
describe conformance requirements. An implementation conforms to this version
only if every externally visible transition preserves the state and financial
invariants defined here.

\section{Scope}

This specification defines:

\begin{itemize}[leftmargin=2em]
    \item private principals and public protocol identities;
    \item directed communication relationships;
    \item signed contact terms;
    \item bonded communication attempts and persistence reserves;
    \item relationship acceptance, rejection, blocking, unblocking, and revocation;
    \item unadmitted reservation cancellation and admitted-bond expiry;
    \item repeated-attempt admission and principal-level continuity;
    \item relationship-solicitation coalescing;
    \item canonical event ordering, idempotency, and atomic settlement; and
    \item native content confidentiality, relationship-graph privacy, record
          lifecycle, minimum ledger semantics, and protocol invariants.
\end{itemize}

Message-body formats, content moderation, inbox ranking, user-interface design,
external funding, withdrawal, foreign exchange, and inter-provider net settlement
are outside the deterministic kernel. Those systems \MUSTNOT{} contradict the
state and settlement consequences defined here.

\section{Protocol Model at a Glance}

\paragraph{Status of this section.}
This section is nonnormative. It gives the reader a map of the interacting state
before the following sections define it precisely. If this overview conflicts
with a later requirement, the later normative requirement controls.

\subsection{One First Contact}

Alice's public protocol identity has not been accepted by Bob's public protocol
identity. Bob's provider quotes $C=0.25$ and $S=4.75$ settlement units. On the
first attempt, $L_0=0$, so Alice reserves $R_0=5.00$. After the encrypted message
is admitted, Bob's relationship decision determines both permission and money:

\begin{center}
\small
\begin{tabularx}{\textwidth}{@{}lYYY@{}}
\toprule
\textbf{Outcome} & \textbf{Bond settlement} & \textbf{Relationship effect} & \textbf{Repeated-attempt effect} \\
\midrule
Accept & Return $C+S$ to Alice. & Establish directed permission for Alice's public identity. & Release active $L_k$ for that identity; preserve principal-level history. \\
Reject & Transfer $C$ to Bob's provider and $S$ to Bob. & Decline this solicitation; permit a later eligible attempt after backoff. & Preserve the advanced level; release each $L_k$ on schedule. \\
Expire & Transfer $C$ to Bob's provider and return $S$ to Alice. & Leave the current relationship state unchanged. & Preserve the advanced level; release each $L_k$ on schedule. \\
\bottomrule
\end{tabularx}
\end{center}

Acceptance and rejection apply to the Alice-to-Bob relationship, not merely to
the message beside which Bob acts. A timely decision therefore settles every
eligible admitted bond for that directed identity pair. A reserved but unadmitted
message is instead cancelled with a full return. The detailed definitions appear
below. The first admitted attempt opens one relationship solicitation; later
attempts join that grouped episode without generating another relationship-
establishment signal.

\subsection{Five Interacting Lifecycles}

The protocol composes five related but distinct lifecycles:

\begin{center}
\small
\begin{tabularx}{\textwidth}{@{}p{0.18\textwidth}YY@{}}
\toprule
\textbf{Lifecycle} & \textbf{States or stages} & \textbf{What changes it} \\
\midrule
Directed relationship & \texttt{unknown}, \texttt{accepted}, \texttt{rejected}, \texttt{revoked}, \texttt{blocked} & Recipient acceptance, rejection, revocation, blocking, or unblocking. \\
Communication attempt & terms issued, value reserved, then admitted or cancelled & Sender authorization followed by durable admission, sender cancellation, or admission timeout. \\
Bond & \texttt{reserved} or \texttt{admitted}, then one terminal state & Unadmitted cancellation, acceptance, rejection, blocking, or expiry. \\
Persistence reserve & \texttt{reserved}, then \texttt{released} exactly once & Immediate release on unadmitted cancellation or identity-scoped acceptance; scheduled release after rejection, blocking, or expiry. \\
Relationship solicitation & inactive, active generation, then closed & First admitted unaccepted attempt opens it; later attempts join it; a relationship decision or exhaustion of decision-open bonds closes it. \\
\bottomrule
\end{tabularx}
\end{center}

The relationship controls future communication. The bond settles one or more
pending attempts. Repeated-attempt state controls when and with how much
additional refundable liquidity another attempt may be admitted. The relationship
solicitation controls whether the recipient receives a new request to establish
the relationship or merely sees another attempt grouped into an existing request.
The ledger records the ownership consequences. No one state is a substitute for
another.

\subsection{Identity and Reference Layers}

A private principal $P_i$ is a provider-local control and anti-abuse subject. A
public protocol identity $I_i$ is the visible persona to which permission is
granted. An opaque relationship handle $H_{ij}$ represents a directed identity
pair inside a limited service domain. Ledger, content, and telemetry systems use
their own scoped references. These layers deliberately do not collapse into one
global user identifier.

\section{Terminology and Actors}

\subsection{Actors}

\begin{description}[leftmargin=2.2em,style=nextline]
    \item[Sender] An authenticated private principal acting through a public
    protocol identity to initiate a communication attempt.
    \item[Recipient] The private principal controlling the destination protocol
    identity and its directed communication relationship.
    \item[Sender provider] The service that authenticates the sender and reserves
    sender funds.
    \item[Recipient provider] The service that publishes contact terms, admits the
    message, records recipient decisions, and delivers to the recipient.
    \item[Canonical journal] The authoritative, append-only ordering of protocol
    and ledger events for a deployment or federated settlement domain.
    \item[Ledger operator] The service maintaining canonical reservations,
    transfers, releases, and balance projections.
\end{description}

One organization \MAY{} perform several roles in a centralized deployment. Their
logical responsibilities remain distinct.

\subsection{Identifiers}

The following notation is used:

\begin{itemize}[leftmargin=2em]
    \item $P_i$: the private principal reference for sender $i$ within its
    authoritative identity service;
    \item $I_i$: sender $i$'s public protocol identity;
    \item $I_j$: recipient $j$'s public protocol identity;
    \item $A_{ij}$: permission state from $I_i$ to $I_j$; and
    \item $(P_i,I_j)$: the private key for repeated-attempt state.
\end{itemize}

The notation $P_i$ denotes the continuing enforcement subject, not a portable or
globally comparable identifier. An implementation \MUSTNOT{} disclose its local
principal reference or underlying identity evidence in ordinary messages, public
relationship records, or inter-provider envelopes. A federated assertion \MAY{}
carry an opaque, scoped reference sufficient to address the correct
repeated-attempt state.

\subsection{Authenticated Message Declarations}

A \emph{declared purpose} is a versioned, machine-readable statement by the
authority that constructs the message envelope about the intended communication
class. It is not a provider determination that the encrypted content truthfully
belongs to that class. The base vocabulary \MUST{} include an explicit
\texttt{unspecified} value; extensions \MUST{} use namespaced, versioned
identifiers rather than uncoordinated free text.

An \emph{origin declaration} identifies the accountable submission mode, such as
human-initiated, agent-drafted and human-approved, autonomous-agent, automated-
system, or legacy/unspecified. Authentication proves which authorized sender or
gateway made the declaration. It does not prove that a human composed particular
words or that the declared workflow occurred. Any stronger attestation is a
separate, purpose-limited extension and \MUSTNOT{} place general legal-identity
evidence in an ordinary message envelope.

A \emph{payload-schema identifier} is an optional namespaced and versioned label
for the structure of encrypted content. It exposes no payload field by itself.
An unknown noncritical schema identifier \MUSTNOT{} invalidate otherwise valid
base delivery. An extension that requires schema understanding \MUST{} be
explicitly advertised and treated as critical before reservation or admission.

The message validity time is the final canonical time at which the message may be
admitted. It is distinct from the recipient decision deadline, application-level
meaning inside encrypted content, and the retention or deletion time of stored
ciphertext. Once admission succeeds, message validity expiry \MUSTNOT{} shorten a
bond decision window, alter settlement, or authorize premature deletion.

\subsection{Relationship-Level Decision}

An acceptance, rejection, or block applies to a directed communicative
relationship, not to one selected message. A user interface \MAY{} place the
action beside a message, but the authenticated command \MUST{} identify the sender
and recipient protocol identities and \MUST{} carry relationship-level
consequences. Rejection declines the active solicitation and preserves the
possibility of a later eligible attempt after backoff. Blocking separately
prohibits future ordinary contact until the recipient removes the block.

\subsection{Relationship Solicitation Episode}

A relationship solicitation episode is the recipient-provider grouping of one or
more admitted attempts from the same unaccepted sender protocol identity to the
same recipient protocol identity. The first admitted attempt in a new episode
produces exactly one logical
relationship-establishment signal. Its delivery
channel may follow recipient settings. Later admitted attempts join the active
episode and do not produce another such signal.

The episode is not a message, bond, permission state, or ledger account. Attempts
inside it remain distinct and are not idempotent with one another. Each retains
its own content, bond, decision deadline, and persistence reserve. Exact replay of
one command remains governed by its idempotency key.

Acceptance, rejection, or blocking closes the active episode. Rejection permits a
later eligible admission after backoff; blocking does not. Silence is not a
relationship command: each bond remains subject to its own expiry. An episode
lapses when no associated admitted bond remains decision-open. A later eligible
admission may open a new episode. A user-interface dismissal without an
authenticated relationship decision has no protocol or settlement consequence.

\section{Canonical Time and Event Ordering}

Every state-changing command \MUST{} have one canonical linearization point in
the journal. The journal \MUST{} assign a strictly ordered position and an
authoritative receipt time. Implementations \MUSTNOT{} use a sender-controlled or
recipient-device clock to determine whether a financial decision was timely.

Commands that appear concurrent are resolved by canonical journal order. A
transaction that reads a relationship or repeated-attempt version \MUST{} either
commit against that version or fail without partial effects.

For an authenticated recipient decision, the decision ingress service \MUST:

\begin{enumerate}[leftmargin=2em]
    \item authenticate the recipient and the targeted relationship;
    \item durably record the command or an equivalent commitment;
    \item assign its authoritative receipt time and ordering position; and
    \item return a signed decision receipt.
\end{enumerate}

Federated deployments \MUST{} define which service is authoritative for decision
ingress and how its signed receipts become canonical. A recipient provider
\MUSTNOT{} earn an expiry settlement by discarding a decision for which it issued
a valid timely receipt.

\section{Protocol State}

\subsection{Principal and Protocol Identity}

A \emph{private principal} is a provider-local authorization and enforcement
subject. It binds authentication, control of protocol identities, recovery state,
and repeated-attempt continuity. It need not denote a verified human or legal
person; provider policy \MAY{} use a person, account, organization, or managed
group as the subject. A private principal record minimally contains a local,
immutable reference, verification status where applicable, authorized public
protocol identities, and recovery state.

A \emph{public protocol identity} is a user-visible, protocol-facing persona
through which a principal sends and receives communication. Its record minimally
contains its identity value, controlling provider, authorized operational keys,
status, and version. Directed relationships and communication permission are
scoped to these protocol identities.

One private principal \MAY{} control multiple public protocol identities.
Acceptance of one identity \MUSTNOT{} imply acceptance of another identity
controlled by the same principal. Conversely, changing protocol identities or
operational keys \MUSTNOT{} by itself reset the private principal's
repeated-attempt state for the same recipient.

\pagebreak
A private principal reference \MUSTNOT{} serve as a public address, relationship
identifier, ledger identifier, telemetry identifier, or globally comparable
inter-provider identifier. Outside the authoritative identity service, the
principal \MUST{} be represented only by purpose-limited, scoped assertions as
defined in the privacy requirements of this specification.

\subsection{Directed Relationship}

The directed relationship state is

\begin{equation}
    A_{ij} \in
    \{\texttt{unknown},\texttt{accepted},\texttt{rejected},\texttt{revoked},
      \texttt{blocked}\}.
\end{equation}

The relationship record \MUST{} contain $I_i$, $I_j$, the current state, a
monotonically increasing version, the most recent canonical event position, and
the transition time.

A bond is required exactly when the relationship is not accepted and no valid
scoped capability authorizes the message:

\begin{equation}
\operatorname{bond\_required}(I_i,I_j)=
\begin{cases}
0, & A_{ij}=\texttt{accepted},\\
1, & A_{ij}\in
     \{\texttt{unknown},\texttt{rejected},\texttt{revoked}\}.
\end{cases}
\end{equation}

The function is intentionally undefined for \texttt{blocked}. A blocked
relationship is ineligible for ordinary contact rather than eligible upon payment.
The recipient provider \MUSTNOT{} issue ordinary contact terms or admit an ordinary
attempt while the block remains active. A block also overrides previously issued
scoped capabilities for the directed pair unless the recipient removes the block
through an authenticated transition.

\subsection{Repeated-Attempt State}

Repeated-attempt state is keyed by $(P_i,I_j)$, not by a replaceable sender
address, operational key, or provider account. It \MUST{} contain:

\begin{itemize}[leftmargin=2em]
    \item attempt level $k$;
    \item earliest next-admission time;
    \item active persistence-reserve references;
    \item a monotonically increasing version; and
    \item the most recent canonical event position and transition time.
\end{itemize}

Changing sender protocol identities, operational keys, or providers \MUSTNOT{}
reset this state. Distinct recipient protocol identities retain distinct state.
Acceptance of one public identity \MUSTNOT{} reset this principal-level state.
While that identity remains accepted, its ordinary messages bypass the state;
other public identities controlled by the principal remain subject to it. Any
decay or forgiveness of old principal-level history \MUST{} be an explicit,
versioned policy transition rather than an implicit effect of acceptance.

\subsection{Contact Terms}

Before reserving an unaccepted attempt, the sender \MUST{} obtain a signed quote
from the recipient provider. A quote \MUST{} contain at least:

\begin{itemize}[leftmargin=2em]
    \item quote identifier and protocol version;
    \item sender and recipient protocol-identity references;
    \item recipient-provider identity;
    \item an opaque repeated-attempt-state reference;
    \item relationship and repeated-attempt versions;
    \item processing charge $C$ and effective collateral $S$;
    \item attempt level $k$, backoff eligibility time, and persistence reserve
          $L_k$;
    \item settlement unit;
    \item admission-window duration, decision-window duration, and persistence-
          reserve release conditions;
    \item quote issue and expiry times;
    \item when eligibility or quoted terms depend on authenticated message
          declarations, a canonical digest binding those declarations;
    \item privacy-profile and retention-policy versions;
    \item policy version; and
    \item provider signature.
\end{itemize}

The recipient provider establishes $C$ and \MAY{} establish minimum collateral
$S^{\min}_{p(j)}$. If the recipient preference is $\widehat S_j$, the effective
collateral is

\begin{equation}
    S=\max\{\widehat S_j,S^{\min}_{p(j)}\}.
\end{equation}

The complete required reservation is

\begin{equation}
    R_k=C+S+L_k.
\end{equation}

Outside the endpoint and routing boundary, the protocol-identity references
\MUST{} be pairwise-scoped and opaque as defined in the privacy requirements
below. The sender \MUST{} be shown $R_k$ before authorization. The quote \MUSTNOT{} add
a sender-provider charge, delivery fee, or clearing fee to protocol settlement.
Once a reservation is created, its incorporated terms \MUST{} remain immutable.
The reservation command materializes an admission deadline from the quoted
admission-window duration and canonical reservation time. Admission materializes
the bond's decision deadline from the quoted decision-window duration and
canonical admission time. Both additions \MUST{} use checked canonical-time
arithmetic.

\subsection{Bond}

A bond \MUST{} contain a unique identifier, quote reference, opaque sender and
recipient references sufficient to resolve the directed relationship, attempt and
message references, $C$, $S$, settlement unit, creation time, admission deadline,
optional canonical admission time and decision deadline, current state, and
current version. The admission time and decision deadline are absent while the
bond remains \texttt{reserved} and become immutable when it becomes
\texttt{admitted}.

Bond state is

\begin{equation}
\begin{split}
\texttt{bond\_state}\in\{&\texttt{reserved},\texttt{admitted},
\texttt{accepted},\texttt{rejected},\\
&\texttt{expired},\texttt{cancelled\_unadmitted}\}.
\end{split}
\end{equation}

The first two states are nonterminal. The remaining states are terminal. A
reserved bond whose admission deadline has passed is cancellation-eligible and
never becomes financially expiry-eligible. An admitted bond whose decision
deadline has passed without a timely relationship decision is logically expiry-
eligible even if an expiry worker has not yet materialized its terminal journal
entries.

Define the unadmitted reserved set

\begin{equation}
\mathcal{R}_{ij}=
\left\{b\;\middle|\;
\begin{array}{l}
b.\text{sender}=I_i,\ b.\text{recipient}=I_j,\\
b.\text{state}=\texttt{reserved}
\end{array}
\right\}.
\end{equation}

For canonical receipt time $t$, define the decision-open admitted set

\begin{equation}
\mathcal{D}_{ij}(t)=
\left\{b\;\middle|\;
\begin{array}{l}
b.\text{sender}=I_i,\ b.\text{recipient}=I_j,\\
b.\text{state}=\texttt{admitted},\\
t\leq b.\text{decision\_deadline}
\end{array}
\right\}.
\end{equation}

Relationship decisions financially settle bonds in $\mathcal{D}_{ij}(t)$.
Bonds in $\mathcal{R}_{ij}$ are cancelled with a full return because no message
was admitted. An admitted bond whose deadline precedes $t$ is logically expired
and is in neither set.

\subsection{Persistence Reserve}

A persistence reserve \MUST{} contain a unique identifier, $(P_i,I_j)$ reference,
the sender public-identity or directed-relationship reference, attempt and quote
references, amount $L_k$, creation time, scheduled release time, state, and
version. Its state is \texttt{reserved} or \texttt{released}.

$L_k$ remains sender property. It \MUSTNOT{} be transferred to a recipient or
provider. Cancellation before admission releases it immediately. Acceptance
releases every active reserve associated with attempts from the accepted public
identity to the recipient identity; it \MUSTNOT{} release a reserve associated
with another public identity controlled by the same principal. Following
rejection, blocking, or expiry, it remains reserved until its quoted release time
and is then returned exactly once.

\subsection{Native Message Envelope and Encrypted Content}

Plaintext message content is not protocol or ledger state. In the native privacy
profile, content comprises the subject, body, inline media, attachments, and any
content-derived search representation. It is encrypted by the sender endpoint and
decrypted only by authorized recipient endpoints.

The native message envelope contains only what is necessary to authenticate and
route the attempt: a versioned opaque message identifier, scoped sender and
recipient references, creation and validity times, an encrypted-content reference,
a declared purpose, an origin declaration, an optional payload-schema identifier,
and the permission, capability, quote, bond, or admission proofs required for that
delivery. The declared purpose and origin declaration may use their explicit
\texttt{unspecified} or legacy values. Every declaration, identifier, validity
time, and proof that affects admission \MUST{} be covered by the authenticated
envelope and cryptographically bound to the encrypted content.

The envelope \MUST{} use the least revealing declaration sufficient for the
advertised capability or recipient-controlled routing rule. It \MUSTNOT{} contain
plaintext content, legal identity evidence, unrelated protocol identities, an
account balance, relationship history, or content-derived classifications that
are not required for the selected delivery path. A payload-schema identifier
describes encrypted structure; it does not move structured payload fields into
the envelope.

\subsection{Relationship Solicitation State}

The recipient provider \MUST{} maintain or deterministically derive solicitation
state for each directed unaccepted relationship. An active episode minimally
contains an opaque episode identifier, the directed relationship handle, a
monotonically increasing generation, associated admitted-attempt references, its
open time, status, version, and the canonical position of its most recent change.

When admission finds no active episode, it \MUST{} open the next generation and
create exactly one authenticated logical relationship-solicitation effect. When
admission finds an active episode, it \MUST{} associate the attempt with that episode and
\MUSTNOT{} create another relationship-establishment effect. It \MAY{} update a
grouped message count, unread state, or other quiet presentation.

The logical effect \MUST{} use scoped references and \MUSTNOT{} contain plaintext
message content, a legal identity, an account balance, or unrelated relationship
history. Recipient settings may choose its delivery channel or suppress an
interruptive push, but the active grouped solicitation remains available to the
recipient.

Acceptance, rejection, and blocking \MUST{} close the active episode. Expiry of
the final decision-open admitted bond in an episode \MUST{} cause the episode to
lapse once the journal can determine that no timely decision applies. Closing or
lapsing an episode moves no value by itself and does not change the settlement of
any bond. Recipient presentation details other than the single-solicitation
constraint remain outside the deterministic kernel.

\section{Financial Semantics}

The protocol bond is

\begin{equation}
    B=C+S,
\end{equation}

and the total attempt reservation is $R_k=B+L_k$. Ledger entries \MUST{} use
double-entry accounting or an equivalent representation that proves value
conservation. Available balance, reserved balance, and owned balance \MUST{} be
distinguished.

The ledger journal \MUST{} be append-only. Corrections \MUST{} be represented by
new compensating entries authorized by a separate rule; historical entries
\MUSTNOT{} be rewritten. No compensation rule \MAY{} silently violate terminal
bond finality.

Protocol settlement defines ownership changes. External deposits, withdrawals,
chargebacks, or bank transfers do not alter the meaning of a protocol transition.
An implementation \MUST{} prevent external funding failure from producing an
unbacked valid reservation.

\section{Commands and Transitions}

The command subsections follow one reading pattern even where the normative text
uses prose rather than identical headings:

\begin{enumerate}[leftmargin=2em]
    \item \textbf{Authority and inputs:} who may issue the command and which
          canonical facts it carries;
    \item \textbf{Preconditions:} versions, time, state, signatures, and funds
          that must be valid before commitment;
    \item \textbf{Atomic outcome:} the complete state and ownership result;
    \item \textbf{Idempotency and concurrency:} how replay or a stale snapshot is
          resolved; and
    \item \textbf{Errors:} the stable failure described in Section 10.
\end{enumerate}

The following table is an index, not a replacement for the normative command
definitions:

\begin{center}
\small
\begin{tabularx}{\textwidth}{@{}p{0.22\textwidth}p{0.27\textwidth}Y@{}}
\toprule
\textbf{Command} & \textbf{Principal gate} & \textbf{Atomic protocol result} \\
\midrule
\texttt{IssueContactTerms} & Current relationship, policy, and attempt state & Signed fixed terms, or a no-bond result for an accepted relationship. \\
\texttt{ReserveAttempt} & Valid current quote and sufficient available value & Reserve value and schedule the short admission deadline. \\
\texttt{AdmitAttempt} & Admission window, durable content, and exact reservation verified & Admit the message, start its decision window, advance attempt state, and open or join one solicitation. \\
\texttt{CancelReserved}\newline\texttt{Attempt} & Authenticated sender cancellation, admission timeout, or pre-admission failure & Return $C+S+L_k$ in full and do not advance attempt state. \\
\texttt{AcceptRelationship} & Authenticated recipient decision at canonical time & Accept the public-identity relationship, settle admitted bonds, and fully cancel unadmitted reservations. \\
\texttt{RejectRelationship} & Authenticated recipient decision while unaccepted & Decline the solicitation, settle decision-open admitted bonds, and fully cancel unadmitted reservations. \\
\texttt{BlockRelationship} & Authenticated recipient decision for a directed pair & Prohibit future ordinary contact, settle admitted bonds, and fully cancel unadmitted reservations. \\
\texttt{UnblockRelationship} & Authenticated recipient controls a blocked relationship & Return to rejected status without accepting, resetting backoff, or reopening bonds. \\
\texttt{ExpireBond} & Deadline closed and no timely decision applies & Expire once, transfer $C$, return $S$, and retain $L_k$ until release. \\
\texttt{ReleasePersistence}\newline\texttt{Reserve} & Scheduled release time reached & Return $L_k$ exactly once. \\
\texttt{RevokeRelationship} & Authenticated recipient controls an accepted relationship & Revoke future bond-free communication without reopening settled bonds. \\
\bottomrule
\end{tabularx}
\end{center}

\subsection{Issue Contact Terms}

\texttt{IssueContactTerms} evaluates the authenticated sender identity, current
relationship, repeated-attempt state, recipient policy, and provider policy.

If any eligibility decision or quoted value depends on declared purpose, origin,
schema requirements, or another authenticated envelope declaration, the request
for terms \MUST{} supply the canonical declaration set and the signed terms
\MUST{} bind its digest. A provider \MUSTNOT{} quote under one declaration set and
admit a materially different set. The reference deployment does not vary bond
economics merely to purchase priority, attention, or a response.

If $A_{ij}=\texttt{accepted}$, the provider \SHOULD{} return a no-bond result and
\MUSTNOT{} require $C$, $S$, or $L_k$.

If $A_{ij}=\texttt{blocked}$, the provider \MUST{} return
\texttt{ContactBlocked} and \MUSTNOT{} issue ordinary contact terms. This rule
takes precedence over a previously issued scoped capability for the pair.

For \texttt{unknown}, \texttt{rejected}, or \texttt{revoked}, the provider
\MUST{} return signed terms consistent with the current relationship and
repeated-attempt versions.

For an initial attempt, $k=0$, $\Delta_0=0$, and $L_0=0$. Protocol policy \MUST{}
define nondecreasing schedules

\begin{equation}
    \Delta_{k+1}\geq\Delta_k,
    \qquad
    L_{k+1}\geq L_k.
\end{equation}

\subsection{Reserve Attempt}

\texttt{ReserveAttempt} consumes a valid quote and an idempotency key. It \MUST:

\begin{enumerate}[leftmargin=2em]
    \item verify the quote signature, version, identities, expiry, and settlement
          unit;
    \item verify that the relationship and repeated-attempt versions still match;
    \item verify that the relationship remains bond-required;
    \item verify that the sender has at least $R_k$ available;
    \item atomically reserve exactly $C+S$ as a bond and $L_k$ as a separate
          persistence reserve;
    \item compute the admission deadline from canonical reservation time and the
          quoted admission-window duration using checked arithmetic;
    \item schedule idempotent cancellation at that deadline; and
    \item return authenticated reservation evidence.
\end{enumerate}

A stale quote \MUST{} fail without retaining value. If acceptance linearized
before this command, the command \MUST{} return \texttt{RelationshipAccepted}; any
provisional hold \MUST{} be released.
If blocking linearized before this command, it \MUST{} return
\texttt{ContactBlocked} and release any provisional hold.

\subsection{Admit Attempt}

\texttt{AdmitAttempt} verifies the message envelope, quote, bond, persistence
reserve, current relationship, repeated-attempt state, and durable encrypted-
content evidence. For an unaccepted relationship it \MUST{} atomically:

\begin{enumerate}[leftmargin=2em]
    \item confirm that canonical admission time does not exceed the envelope's
          message validity time;
    \item verify the integrity and supported criticality of authenticated message
          declarations and any declaration digest fixed by the quote;
    \item confirm that the current time is not earlier than the quoted eligibility
          time;
    \item confirm that canonical admission time is not after the bond's admission
          deadline;
    \item confirm that exactly $C+S+L_k$ is reserved under the quoted terms;
    \item confirm that the encrypted content and a retryable delivery intent can
          be committed durably with this transition;
    \item change the bond from \texttt{reserved} to \texttt{admitted};
    \item record canonical admission time and compute the immutable decision
          deadline from the quoted decision-window duration;
    \item cancel the scheduled unadmitted-reservation timeout and schedule bond
          expiry;
    \item advance the repeated-attempt level;
    \item establish the next eligibility time; and
    \item open a new relationship solicitation episode and enqueue its single
          relationship-establishment effect, or join the existing active episode
          without another such effect; and
    \item record authenticated admission evidence.
\end{enumerate}

Admission is the boundary after which ordinary delivery retries do not reverse
the protocol attempt. Advancement occurs on admission, not later settlement.
Duplicate admission with
the same idempotency key \MUST{} return the original result. A conflicting replay
\MUST{} fail.

If the relationship became accepted after reservation but before admission, the
reserved attempt \MUST{} be cancelled unadmitted and its full $C+S+L_k$ returned.
The message \MAY{} then be admitted without a bond under ordinary accepted-sender
policy. It \MUSTNOT{} advance repeated-attempt state or open an unaccepted-
relationship solicitation episode.

If the relationship became blocked after reservation but before admission, the
block transition has already cancelled the unadmitted reservation and returned
$C+S+L_k$ in full. Admission \MUST{} return \texttt{ContactBlocked}, create no
message or solicitation effect, and perform no additional value movement.

\subsection{Accept Relationship}

\texttt{AcceptRelationship} is authenticated by the recipient and targets
$(I_i,I_j)$. Let $t_a$ be its authoritative receipt time. In one atomic
transition, the command \MUST:

\begin{enumerate}[leftmargin=2em]
    \item set $A_{ij}\leftarrow\texttt{accepted}$ and increment the relationship
          version;
    \item identify every bond in $\mathcal{D}_{ij}(t_a)$ and mark it
          \texttt{accepted};
    \item release exactly $C+S$ for each accepted bond to its sender and cancel
          its scheduled expiry;
    \item identify every bond in $\mathcal{R}_{ij}$, mark it
          \texttt{cancelled\_unadmitted}, return its complete $C+S+L_k$, and
          cancel its admission-timeout operation;
    \item release every other active persistence reserve associated with an
          attempt from $I_i$ to $I_j$, including reserves for earlier expired or
          rejected attempts;
    \item preserve the repeated-attempt level and backoff for $(P_i,I_j)$ and
          preserve every reserve associated with another public identity
          controlled by $P_i$;
    \item close the active relationship solicitation episode and suppress any
          not-yet-delivered relationship-establishment effect for it; and
    \item append the relationship, solicitation, bond, reserve, and ledger events
          as one atomic protocol outcome.
\end{enumerate}

Acceptance settles the relationship, not merely the message beside which the
recipient acted. No bond included in the acceptance transition \MAY{} later be
rejected or expired. No message from $I_i$ admitted after the transition \MAY{}
require a bond while the relationship remains accepted.

A bond whose decision deadline precedes $t_a$ is not in
$\mathcal{D}_{ij}(t_a)$. It \MUST{} settle as expired if it has not already done
so. The acceptance still grants future permission and settles every other
decision-open admitted bond. It \MUSTNOT{} reopen a terminal or logically expired
bond.

Acceptance of $I_i$ \MUSTNOT{} accept another public identity controlled by
$P_i$, reset the principal-level repeated-attempt state, or release persistence
reserves associated with attempts from another public identity.

\subsection{Reject Relationship}

\texttt{RejectRelationship} is authenticated by the recipient and targets
$(I_i,I_j)$. It is invalid while $A_{ij}=\texttt{accepted}$; an accepted
relationship \MUST{} first be revoked. Let $t_r$ be the authoritative receipt time.
In one atomic transition, rejection \MUST:

\begin{enumerate}[leftmargin=2em]
    \item set $A_{ij}\leftarrow\texttt{rejected}$ and increment the relationship
          version;
    \item identify every bond in $\mathcal{D}_{ij}(t_r)$;
    \item mark each identified bond \texttt{rejected};
    \item transfer each identified bond's $C$ to the recipient provider;
    \item transfer each identified bond's $S$ to the recipient;
    \item cancel scheduled expiry operations for those bonds;
    \item identify every bond in $\mathcal{R}_{ij}$, mark it
          \texttt{cancelled\_unadmitted}, return its complete $C+S+L_k$, and
          cancel its admission-timeout operation;
    \item preserve the advanced repeated-attempt level and backoff; and
    \item retain each $L_k$ associated with an admitted bond settled by rejection
          until its quoted release time; and
    \item close the active relationship solicitation episode and suppress any
          not-yet-delivered relationship-establishment effect for it.
\end{enumerate}

Rejection settles every decision-open admitted bond for the relationship at its
canonical ordering point and fully cancels every unadmitted reservation. A
selected message is evidence of which relationship the recipient means; it does
not limit settlement to that message. A bond settled by rejection \MUSTNOT{} later
be accepted or expired.

Rejection is not a permanent contact prohibition. The rejected state remains
bond-required, preserves the repeated-attempt level and backoff, and permits a
later eligible sender attempt to obtain new terms and open a new solicitation.

\subsection{Block Relationship}

\texttt{BlockRelationship} is authenticated by the recipient and targets
$(I_i,I_j)$. It \MAY{} be applied from any relationship state. Let $t_b$ be its
authoritative receipt time. In one atomic transition, blocking \MUST:

\begin{enumerate}[leftmargin=2em]
    \item set $A_{ij}\leftarrow\texttt{blocked}$ and increment the relationship
          version;
    \item identify every bond in $\mathcal{D}_{ij}(t_b)$;
    \item settle each identified bond under the rejection allocation by marking
          it \texttt{rejected}, transferring $C$ to the recipient provider, and
          transferring $S$ to the recipient;
    \item cancel scheduled expiry operations for those bonds;
    \item identify every bond in $\mathcal{R}_{ij}$, mark it
          \texttt{cancelled\_unadmitted}, return its complete $C+S+L_k$, and
          cancel its admission-timeout operation;
    \item preserve the repeated-attempt level, backoff, and scheduled release of
          every $L_k$ associated with an admitted bond settled by blocking;
    \item invalidate scoped capabilities for the directed pair;
    \item close the active relationship solicitation episode and suppress any
          not-yet-delivered relationship-establishment effect for it; and
    \item append the relationship, solicitation, bond, capability, reserve, and
          ledger events as one atomic protocol outcome.
\end{enumerate}

While the relationship remains blocked, the recipient provider \MUSTNOT{} issue
ordinary contact terms, accept a reservation, or admit an ordinary or
capability-authorized attempt for the pair. An exact replay returns the recorded
result without another settlement.

\subsection{Unblock Relationship}

\texttt{UnblockRelationship} is authenticated by the recipient. It \MUST{} change
$A_{ij}$ from \texttt{blocked} to \texttt{rejected} and increment the relationship
version. It \MUSTNOT{} accept the relationship, reset repeated-attempt state,
release persistence reserves early, restore invalidated capabilities, reopen
settled bonds, or create a solicitation. A future sender attempt remains subject
to the preserved backoff and must obtain fresh contact terms.

\subsection{Expire Bond}

Only an admitted bond has a decision deadline $T_b$ and may be expired. A
reserved bond is governed by its admission deadline and
\texttt{CancelReservedAttempt}. A decision is timely for admitted bond $b$ when
its authoritative receipt time is at or before $T_b$. Expiry \MUSTNOT{} become
terminal until the canonical journal can determine that no timely acceptance,
rejection, or blocking decision applies to the bond.

An implementation \SHOULD{} use a bounded processing holdback after $T_b$ before
materializing expiry. The holdback is not an extension of the recipient's decision
window; it allows decisions already durably received by $T_b$ to be sequenced and
observed. The quote or provider policy \MUST{} disclose the earliest time at which
expiry settlement may be considered final.

If a timely acceptance exists, acceptance settlement \MUST{} win. If a timely
rejection exists and no earlier valid acceptance changed the relationship to
accepted, rejection settlement \MUST{} win. A timely block uses the rejection
allocation and then prohibits future contact. Conflicting recipient decisions are
ordered by the canonical journal and relationship-version preconditions.

If neither decision applies, \texttt{ExpireBond} \MUST{} atomically:

\begin{enumerate}[leftmargin=2em]
    \item mark the bond \texttt{expired};
    \item transfer $C$ to the recipient provider;
    \item release $S$ to the sender;
    \item leave the current relationship state unchanged;
    \item preserve the repeated-attempt level and backoff; and
    \item retain the associated $L_k$ until its quoted release time; and
    \item if no associated admitted bond remains decision-open, mark the active
          relationship solicitation episode lapsed without moving value.
\end{enumerate}

The expiry command \MUST{} re-read bond state and relevant relationship decisions
inside its settlement transaction. A stale expiry job encountering an accepted,
rejected, expired, or unadmitted-cancelled bond \MUST{} produce no additional
transfer.

Expiry normally leaves an unknown, rejected, revoked, or blocked relationship in
that state. If a late acceptance has already established prospective permission,
expiry of an older decision-closed bond \MUSTNOT{} undo that acceptance.

\subsection{Cancel Reserved Attempt}

\texttt{CancelReservedAttempt} applies only while the bond remains
\texttt{reserved}. It may be authenticated by the sender as a voluntary
cancellation, by the provider as a verified pre-admission technical failure, or
by a scheduled command whose canonical time is after the admission deadline. In
one atomic transition it \MUST:

\begin{enumerate}[leftmargin=2em]
    \item mark the bond \texttt{cancelled\_unadmitted};
    \item return the full $C+S+L_k$ to the sender;
    \item cancel the admission-timeout operation;
    \item leave repeated-attempt state unchanged; and
    \item create no message, solicitation, decision deadline, or delivery effect.
\end{enumerate}

An exact replay returns the recorded result without another release. A scheduled
timeout that encounters an admitted or terminal bond is a no-op. The sender
\MAY{} cancel before the admission deadline; at or after the deadline the provider
\MUST{} make the same full-return transition if the attempt remains reserved.

After admission, ordinary delivery retry is an outbox concern and
\texttt{CancelReservedAttempt} is invalid. Admission therefore \MUSTNOT{} commit
until encrypted content and a retryable delivery intent are durable. A later
accounting correction requires a separately authorized compensating journal
entry; it does not erase the admitted attempt or roll back repeated-attempt state.

\subsection{Release Persistence Reserve}

At or after a reserve's scheduled release time, an idempotent
\texttt{ReleasePersistenceReserve} command \MUST{} return $L_k$ to the sender if
the reserve remains active. If identity-scoped acceptance or unadmitted
cancellation already released it, the command \MUST{} return the original result
without another ledger entry.

\subsection{Revoke Relationship}

\texttt{RevokeRelationship} is authenticated by the recipient. It \MUST{} change
$A_{ij}$ from \texttt{accepted} to \texttt{revoked} and increment the relationship
version. It \MUSTNOT{} alter previously settled bonds. Future ordinary messages
require a new quote and bond unless a valid scoped capability authorizes them.

\section{Concurrency Rules}

\subsection{Acceptance and Concurrent Reservations}

Acceptance and bond reservation for the same relationship \MUST{} serialize.
A reservation ordered before acceptance is fully cancelled if it remains
unadmitted, or belongs to the decision-open admitted set if admission also
linearized before acceptance. A reservation ordered after acceptance \MUST{} fail
as unnecessary and release any provisional hold. Rejection and blocking use the
same ordering rule, except that decision-open admitted bonds use the rejection
allocation.

\subsection{Admission and Unadmitted Cancellation}

Admission and cancellation for the same reserved attempt \MUST{} serialize. If
admission commits at or before the admission deadline, it cancels the timeout,
starts the decision window, and advances repeated-attempt state. If sender
cancellation commits first, or the attempt remains reserved after the admission
deadline, the full reservation returns and admission cannot later revive it.
There is no state in which an attempt both advances repeated-attempt history and
is cancelled as unadmitted.

\subsection{Acceptance and Expiry}

There is no undefined simultaneous outcome. If acceptance was durably received on
or before a bond's deadline, it \MUST{} take priority even when an expiry worker
began first. If acceptance was received after the deadline, the bond expires and
acceptance applies prospectively and to other bonds whose windows remain open.
Once either settlement commits, retries of the losing command are no-ops with
respect to value.

\subsection{Conflicting Relationship Decisions}

Conflicting acceptance, rejection, blocking, unblocking, or revocation commands
\MUST{} carry an expected relationship version or equivalent compare-and-set
condition. The first command that commits changes the version. A conflicting
stale command \MUST{} fail with \texttt{VersionConflict} and \MUSTNOT{} alter any
bond. A recipient who wishes to decline a newly accepted relationship uses
revocation; a recipient who wishes to prohibit future contact uses blocking.

\subsection{Concurrent Admissions and Solicitation}

Admissions for the same unaccepted directed relationship \MUST{} serialize on the
solicitation generation or use an equivalent uniqueness constraint. Exactly one
admission may change an inactive solicitation to active for a generation and
create its relationship-establishment effect. Every concurrent admission ordered
afterward joins that generation without creating another such effect. A retry of
the winning admission returns its recorded result rather than opening a new
generation.

\subsection{Multiple Pending Bonds}

Relationship settlement \MUST{} enumerate or deterministically derive the complete
decision-open admitted set and the complete unadmitted reserved set at the
transition's journal position. Pagination, sharding, or asynchronous projections
\MUSTNOT{} cause partial relationship settlement or leave an unadmitted
reservation stranded.
Implementations \MAY{} use a transaction manifest or prepare/commit protocol, but
the externally visible result \MUST{} be atomic.

\section{Error Conditions}

An implementation \MUST{} return a stable machine-readable error and \MUSTNOT{}
retain or transfer funds except as explicitly stated below.

\begin{longtable}{p{0.28\textwidth}p{0.62\textwidth}}
\toprule
\textbf{Condition} & \textbf{Required behavior}\\
\midrule
\endhead
\texttt{AuthenticationFailed} & Reject the command without state or ledger
effects.\\
\texttt{InvalidSignature} & Reject the quote, event, assertion, or receipt without
state or ledger effects.\\
\texttt{QuoteExpired} & Refuse new reservation; release any provisional hold.\\
\texttt{QuoteVersionStale} & Refuse reservation or admission and require new
contact terms.\\
\texttt{RelationshipAccepted} & Do not create a bond; release any provisional hold
and continue only under accepted-sender policy.\\
\texttt{ContactBlocked} & Do not issue terms, reserve value, or admit an ordinary
or capability-authorized attempt for the blocked directed pair. Release any
provisional hold.\\
\texttt{BackoffActive} & Refuse admission and report the next eligible time; do not
advance the attempt level.\\
\texttt{InsufficientFunds} & Create no bond or persistence reserve.\\
\texttt{AmountMismatch} & Refuse reservation or admission; no underfunded bond is
valid.\\
\texttt{AdmissionWindowClosed} & If the bond remains reserved, cancel it,
return $C+S+L_k$ in full, and do not admit or advance attempt state.\\
\texttt{ReservedAttempt}\newline\texttt{Cancelled} & Return the recorded full-release result; do
not admit, create a solicitation, or move value again.\\
\texttt{DecisionWindowClosed} & Do not apply the decision to the closed bond; the
relationship command \MAY{} still have prospective effect.\\
\texttt{AlreadyTerminal} & Return the recorded terminal result without another
transfer or release.\\
\texttt{VersionConflict} & Apply no partial relationship or ledger effects; the
caller \MUST{} refresh state before another command.\\
\texttt{Duplicate} & Return the original result when payload and idempotency key
match. A conflicting payload using the same key \MUST{} fail.\\
\texttt{PrivacyScopeViolation} & Refuse an identifier translation, disclosure, or
query that exceeds the caller's authorized relationship and purpose.\\
\texttt{ContentKeyUnavailable} & Do not ask a provider to bypass end-to-end
encryption. Return ciphertext or an endpoint recovery error without changing
settlement.\\
\texttt{MessageValidityClosed} & Refuse admission after the envelope validity
time. Fully release any still-unadmitted reservation, create no solicitation, and
do not advance repeated-attempt state.\\
\texttt{CapabilityScope}\newline\texttt{Mismatch} & Do not consume capability allowance or admit
through that capability. Ordinary bonded admission may proceed only through its
own valid quote and reservation.\\
\texttt{UnsupportedCritical}\newline\texttt{Extension} & Refuse reservation or admission that
depends on an unsupported critical declaration or payload schema; release any
provisional or still-unadmitted hold without advancing attempt state.\\
\texttt{LinkageErased} & Report that expired identifying linkage cannot be
reconstructed; do not restore it from analytics, logs, or an unrelated backup.\\
\bottomrule
\end{longtable}

Provider timeouts, database retries, duplicate network delivery, and process
crashes \MUST{} resolve to one of these deterministic outcomes. An unknown outcome
\MUST{} be queried by idempotency key before the caller retries with a new key.

\section{Launch Extension: Scoped Capabilities and Express Lanes}

The base settlement kernel remains usable without scoped capabilities. The
reference deployment advertises the express-lane capability extension at launch;
an implementation that does not advertise it treats unaccepted ordinary
communication under the bond rules above. When advertised, the following
requirements apply.

A recipient \MAY{} issue a signed communication capability containing scoped
sender and recipient references, a declared-purpose constraint, an optional
origin constraint, validity interval, maximum count or rate, capability version,
acceptable declaration authority, and revocation conditions. Outside the endpoint
and routing boundary, those references \MUST{} be pairwise opaque.

An advertised express-lane capability \MUST{} contain a recipient signature, a
hard validity interval, and an enforceable maximum count or rate. Admission
\MUST{} compare the authenticated declared purpose and any constrained origin
exactly under the capability's versioned matching rules. This comparison enforces
what the sender or gateway declared; it \MUSTNOT{} be represented as provider
inspection or verification of end-to-end encrypted content. A recipient endpoint
\MAY{} report misuse and revoke or block the authority. A native opaque pair
handle identifies the subject but is not authority without the signed grant.

For a legacy sender, a DMARC pass under RFC~9989 is evidence that use of the
RFC 5322 author domain was authorized; it is not evidence that a message is
safe, that its content matches a declared purpose, that a human authored it, or
that a look-alike domain belongs to an expected organization. A legacy ingress
gateway that constructs the authenticated envelope \MUST{} identify itself as the
declaration authority and use a legacy/unspecified origin unless it possesses
separate authenticated evidence. A legacy capability \MUST{} explicitly state
whether such gateway-originated purpose declarations satisfy its scope. A legacy
lane \MUST{} bind separately to a canonical IDNA-ASCII organizational domain and
\MUST{} apply the same versioned domain-to-protocol-identity mapping for grant,
admission, blocking, and revocation. The aligned DKIM or SPF identifier
\MUSTNOT{} silently replace that lane identity.

A message whose authenticated declarations and other constraints match a valid
capability \MAY{} be admitted without a bond while the relationship remains
unaccepted. Capability validation \MUST{} be deterministic, auditable, scoped to
the represented authority, and replay-safe. A declaration mismatch \MUSTNOT{}
consume allowance. A capability
\MUSTNOT{} silently create unrestricted standing acceptance. A capability-
authorized message \MUSTNOT{} open or join a relationship solicitation and
\MUSTNOT{} advance repeated-attempt state.
A \texttt{blocked} relationship overrides a capability for the directed pair;
blocking \MUST{} invalidate that authority until the recipient explicitly issues
new authority after unblocking.

\section{Security, Privacy, and Governance Requirements}

\subsection{Authentication and Replay}

Every quote, reservation assertion, admission receipt, recipient decision,
settlement receipt, and provider assertion \MUST{} be authenticated. Every
state-changing command \MUST{} carry an idempotency key and domain separation
sufficient to prevent reuse in another command type, relationship, or protocol
version.

Operational keys \MUST{} be rotatable without changing $P_i$. Recovery and key
revocation \MUST{} prevent a compromised accepted identity from retaining
indefinite bond-free authority.

\subsection{Native Content Confidentiality}

A conforming native client \MUST{} encrypt message content before transmitting it
to a provider. Only sender-authorized and recipient-authorized endpoints \MAY{} hold
the ordinary content-decryption keys. A provider, ledger operator, content store,
or delivery worker \MUSTNOT{} require possession of those keys as a condition of
native delivery.

The cryptographic binding between the encrypted content and its authenticated
envelope \MUST{} prevent substitution, truncation, reordering, and attachment
replacement. The recipient \MUST{} be able to authenticate the sending protocol
identity and detect an unexpected operational-key change. Implementations
\SHOULD{} provide auditable key transparency or an equivalent mechanism that
makes silent key substitution detectable.

Content storage \MUST{} use ciphertext references rather than embedding content
in relationship or ledger events. Server-side indexes \MUSTNOT{} contain plaintext
subjects, bodies, attachments, or content-derived terms. Search over native
content \MUST{} occur at an authorized endpoint or through a construction that
does not disclose the searched plaintext to the provider.

An SMTP gateway or other legacy bridge \MAY{} necessarily encounter plaintext or
provider-readable content. Such delivery \MUST{} be identified as outside the
native content-confidentiality guarantee. A provider \MUSTNOT{} represent ordinary
TLS transport or provider-managed encryption at rest as native end-to-end
encryption.

\subsection{Pairwise Relationship Privacy}

The logical relationship remains $(I_i,I_j)$, but components that do not route or
display an address \MUST{} represent it by an opaque pairwise handle $H_{ij}$. The
handle \MUST{} be generated randomly or by a secret-key construction with domain
separation and at least 128 bits of effective unpredictability. It \MUSTNOT{} be a
plain, unsalted, or otherwise publicly testable hash of the two protocol
identities.

A handle is scoped to one directed identity pair and one service domain. The
handle used for $I_i$ contacting $I_j$ \MUSTNOT{} reveal the handle used for the
same sender contacting another recipient. The reverse relationship from $I_j$ to
$I_i$ \MUST{} use a distinct handle. Relationship, content, ledger, and
observability systems \SHOULD{} use separately scoped references so that one
global identifier does not join all records without passing through a controlled
translation boundary.

The private repeated-attempt relationship $(P_i,I_j)$ similarly \MUST{} be
represented outside the identity service by a recipient-scoped opaque assertion.
That assertion \MUST{} allow the correct anti-persistence history to survive changes
of sender address, key, or provider, while preventing unrelated recipients from
comparing assertions to discover a common private principal.

A provider migration \MUSTNOT{} export a globally stable principal identifier.
Where continuity across providers is supported, the prior identity authority
\MUST{} instead issue or authorize a migration proof scoped to the recipient and
the repeated-attempt purpose. The receiving provider \MUSTNOT{} be able to use
that proof to discover the principal's relationships with other recipients.

There \MUSTNOT{} be a public or ordinary administrative interface that enumerates
the global relationship graph. A recipient \MAY{} inspect its own visible
relationships, and a sender \MAY{} inspect its own attempts and receipts, but neither
party receives another participant's unrelated edges. Pending, rejected, ignored,
revoked, blocked, and accepted relationships receive the same graph-confidentiality
protection.

An accepted relationship \MAY{} persist as an opaque handle or recipient-issued
capability together with current state, version, and revocation information. The
system \MUSTNOT{} retain the complete formation history merely because current
permission remains active.

\subsection{Compartmentalization and Minimal Disclosure}

The implementation \MUST{} separate identity evidence, relationship state,
encrypted content, financial settlement, and operational telemetry into distinct
access domains. This is a separation of authority as well as storage. Access to
one domain \MUSTNOT{} automatically grant access to the keys or translation data
of another.

Within ordinary operation, the identity service \MAY{} know that a verified
principal controls a public identity, but it need not receive message content.
The relationship service \MAY{} know $H_{ij}$ and its state, but it need not receive
legal identity evidence or plaintext content. The ledger \MAY{} know opaque account,
bond, and settlement references together with amounts, but it need not receive a
subject, content reference, public address pair, or private principal identifier.
The content store \MAY{} know a ciphertext object and a scoped retrieval reference,
but it need not receive a bond outcome or readable endpoint identities.

The delivery boundary necessarily learns enough to route to $I_j$ and to present
$I_i$ to the recipient. It \MUST{} disclose no additional principal evidence,
unrelated identities, balance information, or relationship history. A recipient
provider \MUST{} verify the applicable reservation or capability, but it \MUSTNOT{}
receive the sender's overall balance or unrelated settlement activity. External
financial rails \SHOULD{} see account funding and withdrawal rather than
individual message events. Internal transfers \SHOULD{} be netted or aggregated
before external settlement when correctness and applicable obligations permit.

Translation among scoped references is itself sensitive processing. Every
translation \MUST{} be authorized for a specific protocol purpose and recorded in
an access audit that does not reproduce the translated graph. Bulk translation,
graph export, and use of relationship data for advertising, sale, or unrelated
profiling are outside the protocol purpose and \MUSTNOT{} be enabled by default.

\subsection{Metadata Leakage and Traffic Analysis}

Pairwise handles do not conceal every observable fact. Timing, network addresses,
message sizes, uncommon bond amounts, and synchronized journal events can permit
correlation without a direct identifier. Implementations \MUST{} document which
of these signals each actor can observe and \MUST{} retain them only while needed
for routing, security, settlement, or a disclosed obligation.

Where it does not compromise correctness, an implementation \SHOULD{} reduce
correlation through coarse timestamps, batching, settlement netting, content
padding, separation of network and account logs, and short retention of delivery
telemetry. These protections \MUSTNOT{} alter quoted amounts, obscure a settlement
from its owner, or make bond finality ambiguous.

\subsection{Record Lifecycle and Retention}

Every stored field \MUST{} belong to a documented retention class stating its
purpose, access domain, controlling key, operational lifetime, dispute lifetime,
any separately applicable required-retention period, and deletion method. A vague
interest in future analytics is not a protocol purpose for indefinite retention.

While an attempt is pending, the protocol \MAY{} retain the scoped linkage required
to authenticate admission, prevent replay, apply repeated-attempt policy, receive
a decision, and settle the bond. Once the attempt is terminal and its defined
dispute period ends, the implementation \MUST{} remove fields and indexes that are
not needed for continuing permission, accounting, fraud handling, or an applicable
retention obligation. A durable accounting record \SHOULD{} contain only an
opaque settlement reference, amount, settlement unit, terminal state, canonical
time, and the minimum ownership evidence required for audit.

Continuing acceptance does not justify continuing retention of individual contact
attempts. After acceptance, the live relationship \MAY{} retain $H_{ij}$, current
state, version, capability or key references, and revocation material. Formation
messages, quote details, delivery telemetry, and superseded relationship events
\MUST{} follow their own shorter retention schedules unless independently needed.

Records retained solely for a legal, accounting, fraud, or dispute obligation
\MUST{} be segregated from online delivery and product analytics. They \MUST{} use
separate encryption keys, narrowly authorized access, an explicit expiration
condition, and an audit trail. A hold on one record or account \MUSTNOT{} silently
extend retention for unrelated users or relationships.

Deletion \MUST{} cover primary storage, replicas, caches, search indexes,
telemetry, derived datasets, and backups according to a documented propagation
schedule. If immediate removal from an immutable backup is impracticable, the
backup \MUST{} remain inaccessible to ordinary processing, expire within a fixed
period, and lose access to deleted record keys no later than that period.

\subsection{Cryptographic Erasure and Consent-Locked Archives}

Encryption is not erasure while the operator retains both a usable key and enough
linkage to identify the parties. When a retention period ends, the implementation
\SHOULD{} delete the unnecessary record and destroy every operator-held key that
could recover its identifying linkage. A claim of cryptographic erasure is valid
only when replicas, wrapped key copies, recovery escrow, and backups can no longer
restore that linkage.

The protocol \MAY{} support a consent-locked historical archive in which
reconstruction requires independent sender and recipient key shares. Such an
archive \MUST{} be separate from live permission and settlement state. Refusal or
loss of a share \MUSTNOT{} prevent the protocol from delivering under current
permission, applying revocation, enforcing backoff, or completing a pending
settlement.

If an implementation offers two-party reconstruction, it \MUST{} state precisely
which fields it covers and whether any separately retained compliance record can
identify the transaction without both shares. It \MUSTNOT{} promise two-party-only
reconstruction for information that the operator knowingly retains under an
independent decryption key. Required records that cannot be consent-locked \MUST{} be
minimized and isolated under the preceding retention rules.

No deletion or threshold scheme can erase a message or receipt independently kept
by a sender, recipient, endpoint backup, or external provider. The protocol's
guarantee concerns records controlled by conforming providers and \MUST{} describe
that boundary accurately.

\subsection{Operational Logs and Analytics}

Logs \MUSTNOT{} contain plaintext content, private identity evidence, raw key
material, full account balances, or a reusable global relationship identifier.
Debugging and tracing identifiers \MUST{} be scoped and rotated so that an
observability system does not become an accidental relationship graph.

Analytics \SHOULD{} operate on aggregates produced after the shortest useful
window. A dataset is not anonymous merely because names and addresses have been
removed; if scoped handles, rare amounts, precise timestamps, or auxiliary data
permit practical reidentification, it remains protected relationship data and
\MUST{} follow the corresponding access and retention rules.

\subsection{Provider and Recipient Incentives}

A provider receiving $C$ has an economic interest in rejection or expiry. A
recipient receiving $S$ may have an interest in attracting and rejecting bonded
contact. Deployments \MUST{} make settlement auditable and \MUST{} define fraud,
collusion, reversal, suspension, and dispute policies outside the deterministic
kernel.

Recipient providers \MUST{} present acceptance as an effective available action,
\MUST{} provide signed decision receipts, and \MUSTNOT{} suppress a timely
acceptance while claiming expiry. Pricing policy \SHOULD{} define bounds or review
mechanisms for $C$ and provider-imposed minimum $S$.

\subsection{Delivery Meaning}

A valid bond authorizes a protocol-valid contact attempt. It does not guarantee
inbox prominence, human attention, a response, or exemption from malware,
harassment, unlawful-content, or safety controls. Providers \MUST{} distinguish
financial validity from content and transport policy.

Those controls do not authorize provider decryption of native content. They \MAY{}
operate on the authenticated envelope, endpoint evaluation, recipient reports,
rate and relationship state, or content a participant deliberately discloses for
investigation. Any disclosed plaintext \MUST{} be scoped to that purpose and its
own retention class.

\subsection{Policy Authority and Isolation}

Implementations \MUST{} distinguish three policy authorities. \emph{Protocol
policy} is a versioned, immutable input to deterministic quotation and settlement
transitions. Terms derived from it remain fixed as specified above. \emph{Admission
policy} runs at an authenticated provider or capability boundary and may use only
the envelope, permission proofs, rate and relationship state, safety evidence, and
other data deliberately available to that boundary. It \MUSTNOT{} decrypt native
content or create unrecorded settlement effects. \emph{Presentation policy} runs
for the recipient, preferably at an authorized endpoint, and may control inbox
placement, notifications, digests, ranking, and local content processing. It
\MUSTNOT{} change protocol validity, relationship state, bond ownership, or
canonical history.

Neither admission nor presentation policy \MAY{} impose a protocol bond on
ordinary communication while the directed relationship remains accepted.
Narrowing that permission requires an authenticated relationship transition or a
separately versioned recipient-granted capability model. A policy implementation
may be assisted by probabilistic or learned software, but such output is not by
itself an authenticated protocol command. Any policy result that affects protocol
admission \MUST{} reduce to a versioned, auditable decision over authorized inputs.

Sender-facing explanations \SHOULD{} disclose the actionable constraint needed to
correct a request while withholding unrelated recipient rules, private signals,
and details that would enable enumeration or adaptive probing. Internal audit
records receive their own scoped references and retention class.

\section{Interoperability Requirements}

A federated implementation \MUST{} provide versioned representations for:

\begin{itemize}[leftmargin=2em]
    \item provider and protocol identities;
    \item contact-term quotes;
    \item message envelopes and content references;
    \item declared-purpose, origin-declaration, payload-schema, declaration-
          authority, and extension-criticality representations;
    \item bond and persistence-reserve commitments;
    \item repeated-attempt assertions;
    \item admission and unadmitted-cancellation receipts;
    \item authenticated relationship decisions;
    \item settlement receipts; and
    \item provider-to-provider clearing obligations.
\end{itemize}

Federation \MUST{} preserve the native privacy profile across provider boundaries.
A provider assertion identifying $P_i$ for repeated-attempt enforcement \MUST{} be
scoped to $I_j$ or its recipient-provider privacy domain and \MUSTNOT{} provide a
stable value that other recipients can compare. Settlement and admission receipts
\MUST{} use the minimum scoped references needed by their verifiers. A downgrade
to SMTP, provider-readable content, or a less private metadata profile \MUST{} be
shown to the sender before transmission and represented in the message envelope.

The sender provider \MUST{} prove reservation of the quoted $C+S+L_k$. The
recipient provider \MUST{} verify that proof and the current quote state before
admission. Inter-provider clearing transfers obligations created by $C$ and $S$;
$L_k$ remains sender property and \MUSTNOT{} enter provider revenue or net clearing.

Independent implementations \MAY{} use different storage engines, internal ledger
software, and external financial rails, provided that identical canonical inputs
produce equivalent protocol and ownership transitions.

\section{Required Invariants}

Every conforming implementation \MUST{} preserve all of the following:

\begin{enumerate}[leftmargin=2em]
    \item \textbf{Bond finality.} Each bond reaches at most one terminal state.
    \item \textbf{Reserve finality.} Each persistence reserve is released exactly
          once and is never transferred to a recipient or provider.
    \item \textbf{Value conservation.} Every transition conserves ledger value.
    \item \textbf{Exact funding.} An unaccepted ordinary attempt is not valid
          unless exactly $C+S+L_k$ is reserved under valid fixed terms.
    \item \textbf{Relationship-wide acceptance.} Acceptance atomically establishes
          permission for $(I_i,I_j)$, refunds every decision-open admitted bond,
          fully cancels every unadmitted reservation for the pair, and releases
          active persistence reserves associated with attempts from $I_i$ to
          $I_j$.
    \item \textbf{No later rejection.} A bond or message included in acceptance
          cannot subsequently be rejected or expired.
    \item \textbf{Relationship-wide rejection.} Rejection atomically settles every
          decision-open admitted bond for $(I_i,I_j)$ according to the rejection
          allocation and fully cancels every unadmitted reservation for the pair.
    \item \textbf{Retryable rejection.} Rejection preserves repeated-attempt
          backoff and permits a later eligible bonded attempt; it is not a block.
    \item \textbf{Recipient blocking.} Blocking settles every decision-open
          admitted bond under the rejection allocation, fully cancels every
          unadmitted reservation, and prevents terms, reservation, capability use,
          and admission for the directed pair until authenticated unblocking.
    \item \textbf{Unadmitted full return.} Cancellation before admission returns
          exactly $C+S+L_k$, creates no solicitation or decision deadline, and
          does not advance repeated-attempt state.
    \item \textbf{Timely-decision priority.} A decision durably received by a
          bond's deadline is observed before expiry becomes final.
    \item \textbf{No retroactive reopening.} A later acceptance may grant future
          permission but cannot reopen a terminal or logically expired bond.
    \item \textbf{Atomic attempt advancement.} Admission and repeated-attempt
          advancement occur atomically.
    \item \textbf{Principal continuity.} Changing a sender address, key, or provider
          does not reset $(P_i,I_j)$ attempt state.
    \item \textbf{Identity-scoped permission.} Acceptance of one public identity
          does not accept other identities controlled by the same principal,
          reset principal-level attempt state, or release their persistence
          reserves.
    \item \textbf{Principal non-exportability.} A private principal reference is
          not a public address or globally comparable inter-provider identifier;
          other domains receive only purpose-limited, scoped assertions.
    \item \textbf{Free accepted communication.} Accepted ordinary communication
          requires no protocol bond and incurs no protocol message fee.
    \item \textbf{Replay safety.} Duplicate commands, events, receipts, and expiry
          jobs cannot duplicate value movement or state transition.
    \item \textbf{Declaration integrity.} Every purpose, origin, schema,
          declaration-authority, validity, and capability reference used for
          admission is authenticated and bound to the admitted encrypted content.
    \item \textbf{Fixed declaration-dependent terms.} When contact terms depend on
          message declarations, reservation and admission use exactly the
          declaration set fixed by those signed terms.
    \item \textbf{Capability-scope enforcement.} A capability admits only messages
          whose authenticated declarations match its constraints, and mismatch
          consumes no allowance or settlement value.
    \item \textbf{Message-validity isolation.} Envelope validity gates admission
          but cannot shorten a decision window or rewrite settlement after
          admission.
    \item \textbf{Atomic visibility.} No observer may see permission change without
          its required bond and reserve settlements, or vice versa.
    \item \textbf{Native content blindness.} Providers cannot decrypt a native
          message's subject, body, inline media, or attachments with
          provider-controlled keys.
    \item \textbf{Pairwise unlinkability.} An opaque handle or principal assertion
          disclosed for one recipient cannot be used by another recipient to
          recognize the same relationship or private principal.
    \item \textbf{Graph non-enumerability.} No public or ordinary administrative
          operation exposes the global relationship graph or contact-attempt
          history.
    \item \textbf{Data compartmentalization.} Routine access to content, identity,
          relationships, the ledger, or telemetry does not automatically grant
          access to the other privacy domains or their translation keys.
    \item \textbf{Policy-authority isolation.} Presentation and admission policy
          cannot mutate settlement or canonical relationship history except
          through an independently authenticated protocol command.
    \item \textbf{Bounded retention.} Every identifying field has a declared
          purpose, retention class, expiration condition, and deletion path across
          primary and derived storage.
    \item \textbf{Honest erasure.} A record described as cryptographically erased
          cannot be reconstructed from operator-held keys, key escrow, replicas,
          logs, or backups.
    \item \textbf{Rail minimization.} External funding and withdrawal systems do
          not receive message content or a per-message relationship history from
          the protocol.
    \item \textbf{Solicitation coalescing.} One active solicitation generation for
          a directed unaccepted relationship produces exactly one logical relationship-
          establishment effect; later admitted attempts join it without another
          such effect.
    \item \textbf{Attempt distinctness.} Solicitation coalescing does not merge
          bonds, deadlines, content, persistence reserves, or attempt advancement,
          and closing an episode does not itself move value.
\end{enumerate}

\section{Reference Transition Summary}

For all bonds in the applicable decision-open admitted set, the relationship
outcomes are:

\begin{equation}
\begin{array}{lll}
\text{Accept relationship} &:& C+S\rightarrow\text{sender},\\[3pt]
\text{Reject relationship} &:& C\rightarrow\text{recipient provider},\quad
                                  S\rightarrow\text{recipient},\\[3pt]
\text{Block relationship} &:& C\rightarrow\text{recipient provider},\quad
                                  S\rightarrow\text{recipient},\\[3pt]
\text{Expire bond} &:& C\rightarrow\text{recipient provider},\quad
                       S\rightarrow\text{sender}.
\end{array}
\end{equation}

Every unadmitted reservation affected by cancellation or a relationship decision
returns $C+S+L_k$ in full. Acceptance releases every other active $L_k$ associated
with attempts from the accepted $I_i$ to $I_j$ and preserves principal-level
attempt state. Rejection, blocking, and expiry retain the $L_k$ associated with
admitted attempts until their scheduled release.
The first admitted attempt in a solicitation generation creates exactly one
logical relationship-establishment effect. Later attempts join that generation
without a new solicitation. Acceptance, rejection, or blocking closes it;
absence of a decision leaves each admitted bond to expire separately and the episode
lapses when none remains decision-open. Rejection permits another eligible
attempt after backoff; blocking prohibits one until the recipient unblocks.

\section{Implementation Sequence}

A reference implementation \SHOULD{} proceed in the following order:

\begin{enumerate}[leftmargin=2em]
    \item implement the in-memory deterministic relationship, bond, reserve, and
          repeated-attempt transitions;
    \item test full event sequences, concurrent decisions, and the required
          invariants with unit and property-based tests;
    \item add the append-only double-entry journal and idempotent settlement;
    \item add signed identities, authenticated minimal message declarations,
          endpoint content encryption, key transparency, quotes, admission
          receipts, and decision receipts;
    \item introduce pairwise handles and enforce the identity, relationship,
          content, ledger, and telemetry access boundaries;
    \item persist journal streams and materialized projections with retention
          classes, key separation, and tested cryptographic erasure;
    \item add delivery, encrypted content storage, solicitation coalescing, and
          expiry processing;
    \item add express-lane scoped capabilities without changing base settlement,
          relationship, or repeated-attempt semantics;
    \item add visibly downgraded SMTP compatibility gateways; and
    \item add privacy-preserving federation and clearing after centralized
          semantics are stable.
\end{enumerate}

The first milestone is complete only when it can deterministically execute quote,
reservation, admission, relationship-wide acceptance and rejection, expiry,
unadmitted cancellation and full release, deferred reserve release,
principal-level attempt-state preservation, revocation, blocking, unblocking, and
the opening, joining, and closing of solicitation episodes while preserving every
invariant in this document. Scoped capabilities are not required inside this
settlement milestone. A production milestone additionally requires native end-to-end content
encryption, pairwise relationship references, compartmentalized access, and
verified deletion across replicas and backups.

\end{document}
